\documentclass[10pt, a4paper]{article}

%%%%%%%%%%%%%% CONFIG %%%%%%%%%%%%%%%
%\usepackage[ngerman]{babel}
%\usepackage[english]{babel}
\newcommand{\GroupNumber}{90} 
\newcommand{\AssignmentNumber}{2}
\newcommand{\Semester}{SS 2025}
%%%%%%%%%%%%%%%%%%%%%%%%%%%%%%%%%%%%%

\usepackage{amssymb, amsmath, amsthm}
\usepackage[utf8]{inputenc}
\usepackage{fancyhdr}
\usepackage{xcolor}
\usepackage{enumitem}
\usepackage[parfill]{parskip}
\usepackage{fullpage}
\usepackage[outline]{contour}
\usepackage{ulem}
\usepackage{listings}



\pagestyle{fancy}
\setlength{\headheight}{12.4pt}
\setlength{\headsep}{1.5\headheight}
\fancypagestyle{plain}{
  \fancyfoot{}
  \lhead{Assignment \AssignmentNumber}
  \rhead{Page \thepage}
}
\normalem

\author{Khaled Elnaggar 12033937,\\ Ilma Bajramovic, \\Nhat Minh Hoang 12029170}
\date{}
\title{
  Softwareparadigmen \Semester,
  Übungsblatt \AssignmentNumber \\
  Gruppe \GroupNumber
}

\newcounter{bsp}
\setcounter{bsp}{0}
\newcommand{\task}{\stepcounter{bsp}\subsection*{Beispiel \arabic{bsp}}}
\newcommand{\subtask}[2]{\textbf{#1} #2}

\newcommand{\ul}[1]{%
  \begingroup%
  \renewcommand{\ULdepth}{2.1pt}%
  \contourlength{0.35pt}%
  \uline{\phantom{#1}}\llap{\contour{white}{#1}}%
  \endgroup%
}
\newcommand{\Prog}[1]{\texttt{\ul{#1}}}

\newcommand{\Interp}[2][]{I(\delta,\omega#1,\Prog{#2})}
\newcommand{\ValueOf}[2][]{\omega#1(\Prog{#2})}
\newcommand{\equMod}[2]{\omega#1' \sim_{#2} \omega#1}
\newcommand{\equModAss}[2]{\equMod{#1}{#2}, \omega#1'(\texttt{#2}) = }

\newcommand{\quantorOmega}[2]{\ \omega#1',\ \equMod{#1}{#2}:\ }

\newcommand{\ifThen}[1]{\Prog{if\ #1\ then}\\}
\newcommand{\elseThen}{\Prog{else}\\}
\newcommand{\elseIfThen}[1]{\Prog{else\ if\ #1\ then}\\}
\newcommand{\ifThenElse}[3]{\Prog{if #1 then } #2 \Prog{ else } #3}

\newcommand{\caseIf}[2]{#1 & \text{if } #2 \\}
\newcommand{\caseElseIf}[2]{#1 & \text{else if } #2 \\}
\newcommand{\caseElse}[1]{#1 & \text{otherwise}}


%%%%%%%%%%%%%%%%%%%%%%%%%%%%%%%%%%%%%%%%%%%%%%%%%%#

\begin{document}
\pagestyle{plain}
\maketitle

\task
\textbf{1)}Below is the implementation of the function \Prog{implies(a,b)}.
\begin{align*}
  \ & \delta\Prog{implies(a,b)} = \\
 1\ & \quad \ifThen{isFalse?(a)} 
 2\ & \quad \quad \Prog{T} \\
 3\ & \quad \elseThen 
 4\ & \quad \quad \Prog{b}
\end{align*}

The interpretation of \Prog{{implies(a,b)}} under the given environment $\omega$ and $\delta$ proceeds as follows.

\begin{align*}
\ &  I(\delta, \omega, \Prog{implies(a,b)}) \\
\ &= I(\delta, \omega^1, \Prog{if isFalse?(a) then } li_2 \Prog{ else } li_4) \\
\ &\quad \text{with } \omega^1: \\
\ &\quad \omega1(\Prog{a}) = I(\delta, \omega, \Prog{a}) = \omega(\Prog{a}) = T \\
\ &\quad \omega1(\Prog{b}) = I(\delta, \omega, \Prog{b}) = \omega(\Prog{b}) = T \\
\ &= I(\delta, \omega1, \Prog{b}) \\
\ &\quad \text{since } I(\delta, \omega^1, \Prog{isFalse(a)}) = \text{F} \\
\ &= \omega1(\Prog{b}) \\
\ &= T
\end{align*}

\textbf{2)} Below is the implementation of the function \Prog{any(l)}.
\begin{align*}
  \ & \delta\Prog{any(l)} = \\
 1\ & \quad \ifThen{empty?(l)} 
 2\ & \quad \quad \Prog{F} \\
 3\ & \quad \elseThen 
 4\ & \quad \quad \Prog{or(first(l), any(rest(l)))}
\end{align*}

The interpretation of \Prog{any(l)} under the given environment $\omega$ and $\delta$ proceeds as follows:

\begin{align*}
\ & I(\delta, \omega, \Prog{any}(l)) \\
\ &= I(\delta, \omega^1, \Prog{if empty?(l)} \Prog{ then } la_2 \Prog{ else } la_4) \\
\ &\quad \textit{with } \omega^1: \\
\ &\quad \omega1(\Prog{l}) = I(\delta, \omega, \Prog{l}) = \omega(\Prog{l}) = [\text{F}, \text{T}] \\
\ &= I(\delta, \omega1, \Prog{or(first(l),} \Prog{any(rest(l)}))) \\
\ &\quad \textit{since } I(\delta, \omega1, \Prog{empty(l)}) = \text{F} \\
\ &= \textit{or}(I(\delta, \omega1, \Prog{first(l)}), I(\delta, \omega1, \Prog{any(rest(l)}))) \\
\ &= \textit{or}(\textit{first}(I(\delta, \omega1, \Prog{l})), I(\delta, \omega1, \Prog{any(rest(l)}))) \\
\ &= \textit{or}(\textit{first}(\omega1(l)), I(\delta, \omega1, \Prog{any(rest(l)}))) \\
\ &= \textit{or}(\textit{first}([\text{F}, \text{T}]), I(\delta, \omega1, \Prog{any(rest(l)}))) \\
\ &= \textit{or}(\text{F}, I(\delta, \omega1, \Prog{any(rest(l)}))) \quad \text{(see NR 1)} \\
\ &= \textit{or}(\text{F}, \text{T}) \\
\ &= \text{T}
\end{align*}

\begin{align*}
\ & \textbf{NR 1: } I(\delta, \omega1, \Prog{any(rest(l)}))\\
\ &= I(\delta, \omega^2, \Prog{if empty?(l)} \Prog{ then } la_2 \Prog{ else } la_4) \\
\ &\quad \textit{with } \omega^2: \\
\ &\quad \omega2(\Prog{l}) = I(\delta, \omega1, \Prog{rest(l)}) = \textit{rest}(I(\delta, \omega1, \Prog{l})) = \textit{rest}([\text{F}, \text{T}]) = [\text{T}] \\
\ &= I(\delta, \omega2, \Prog{or(first(l),} \Prog{any(rest(l)}))) \\
\ &\quad \text{since } I(\delta, \omega2, \Prog{empty(l)}) = \text{F} \\
\ &= \textit{or}(I(\delta, \omega2, \Prog{first(l)}), I(\delta, \omega2, \Prog{any(rest(l)}))) \\
\ &= \textit{or}(\textit{first}(I(\delta, \omega2, \Prog{l})), I(\delta, \omega2, \Prog{any(rest(l)}))) \\
\ &= \textit{or}(\textit{first}(\omega2(\Prog{l})), I(\delta, \omega2, \Prog{any(rest(l)}))) \\
\ &= \textit{or}(\textit{first}([\text{T}]), I(\delta, \omega2, \Prog{any(rest(l)}))) \\
\ &= \textit{or}(\text{T}, I(\delta, \omega2, \Prog{any(rest(l)}))) \quad \text{(see NR 2)} \\
\ &= \textit{or}(\text{T}, \text{F}) \\
\ &= \text{T}
\end{align*}


\begin{align*}
\ & \textbf{NR 2: } I(\delta, \omega2, \Prog{any(rest(l)})) \\
\ &= I(\delta, \omega^3, \Prog{if empty?(l) then } la_2 \Prog{else} la_4) \\
\ &\quad \textit{with } \omega^3: \\
\ &\quad \omega3(\Prog{l}) = I(\delta, \omega2, \Prog{rest(l)}) = \textit{rest}(I(\delta, \omega2, \Prog{l})) = \textit{rest}([\text{T}]) = [ ] \\
\ &= I(\delta, \omega3, \text{F}) \\
\ &= \text{F} \\
\ &\quad \textit{since } I(\delta, \omega3, \Prog{empty(l)}) = \text{T}
\end{align*}



\task
Induction proof of the \texttt{pow(a,b)} function: \\
\begin{enumerate}
    \item  Base case: $\omega(y) = 0 $

\begin{align*}
& I(\delta, \omega, \underline{\texttt{pow}(x, y))} \\
    = \ & I(\delta, \omega, \ifThenElse{eq?(y,0)}{1}{mul(x,pow(x,minus(y,1)))} \\
= \ & I(\delta, \omega, 1) \\
& \quad \textit{since } I(\delta, \omega, \underline{\texttt{eq?}(y, 0))} \implies eq?(\omega(y), 0) \implies eq?(0, 0) \xrightarrow{} \text{true} \\
= \ & 1 \\
= \ & \omega(x)^0 \\
= \ & \omega(x)^{\omega(y)}
\end{align*}

\item Hypothesis to prove: 
\[
\forall \omega \in \text{ENV} : 
I(\delta, \omega, \text{pow}(x, y)) = \omega(x)^{\omega(y)} 
\quad \text{if} \quad \omega(x) \geq 0, \omega(y) \geq 0
\]

\item Inductive step: $\omega(y) = k + 1$


\begin{align*}
& I(\delta, \omega, \underline{\texttt{pow(x,y)}}\\
= \ & I(\delta, \omega, \ifThenElse{eq?(y,0}{1}{\texttt{mul(x,pow(x,minus(y,1)))}} \\
= \ & I(\delta, \omega, \underline{\texttt{mul(x, pow(x, minus(y, 1)))}}) \\
& \quad \textit{since } I(\delta, \omega, \underline{\texttt{eq?(y, 0)}}) \implies eq?(\omega(y), 0) \implies eq?(k+1, 0) \xrightarrow{} \textit{false} \\
= \ & I(\delta, \omega', \underline{\texttt{mul(a, pow(a, b))}} \\
&\quad \omega'(a) = I(\delta, \omega, \underline{\texttt{{x}}}) \implies \omega(x) \\
&\quad \omega'(b) = I(\delta, \omega, \underline{minus(\omega(y), \omega(1))} \implies \omega(y) - \omega(1) = k + 1 -1 = k \\
= \ & I(\delta, \omega'', \underline{\texttt{mul(a, b)}}) \\
&\quad \omega''(a) = I(\delta, \omega', \underline{\texttt{a}}) \implies \omega'(a) \implies \omega(x) \\
&\quad \omega''(b) = I(\delta, \omega', \underline{\texttt{pow(a,b)}}) \implies pow(\omega'(a), \omega'(b)) = \omega'(a)^{\omega'(b)} = \omega(x)^k = \omega(x)^{\omega(y)}
= &\omega(x)*\omega(x)^k \\
=&\omega(x)^{k+1} \\
=&\omega(x)^{\omega(y)}
\end{align*}
\end{enumerate}

Induction proof of the \texttt{div(a,b)} function:
\begin{enumerate}
    \item Base case: $\omega(x)=0$

\begin{align*}
& I(\delta, \omega, \underline{\texttt{div}(x, y))} \\
    = \ & I(\delta, \omega, \ifThenElse{lt?(x,y)}{0}{\texttt{plus(div(minus(x,y),y),1)}} \\
= \ & I(\delta, \omega, 0) \\
& \quad \textit{since } I(\delta, \omega, \underline{\texttt{lt?}(x, y))} \implies lt?(\omega(x), \omega(y)) \implies lt?(0, \omega(y)) \xrightarrow{} \textit{true} \\
= \ & 0 < \omega(y) \xrightarrow{} \textit{true, since $\omega(y)>0$} \\
\end{align*}   
\item Hypothesis to prove: 
\[
\forall \omega \in \mathrm{ENV}:\ I(\delta, \omega, \mathrm{div}(\mathbf{x}, \mathbf{y})) = \lfloor \omega(\mathbf{x}) \div \omega(\mathbf{y}) \rfloor \quad \text{if } \omega(\mathbf{x}) \geq 0,\ \omega(\mathbf{y}) > 0
\]

\item Inductive step $\omega(x)=k+1$
\begin{align*}
  & I(\delta, \omega, \text{div}(x, y)) \\
    = \ & I(\delta, \omega, \ifThenElse{lt?(x,y)}{0}\texttt{plus(div(minus(x,y),y),1}\\ 
    = \ & I(\delta, \omega, \underline{\texttt{plus(div(minus(x,y),y),1}}\\
    & \quad \textit{Case 1: $ \omega(x) < \omega(y)$ } I(\delta, \omega, \underline{\texttt{lt?(x,y)}}) \implies lt?(\omega(x), \omega(y)) \implies k + 1 < \omega(y) \implies 0 \xrightarrow{} true \\
    & \quad \textit{Case 2: $\omega(x) \geq \omega(y)$ } I(\delta, \omega, \underline{\texttt{lt?(x,y)}}) \implies lt?(\omega(x), \omega(y)) \implies k+1>\omega(y) \xrightarrow{} false \\
= & I(\delta, \omega, \underline{\texttt{plus(div(minus(x,y),y),1)}}) \\
= & I(\delta, \omega', \underline{\texttt{plus(div(a,b),1)}} \\
& \quad \omega'(a) = I(\delta, \omega, \underline{\texttt{minus(x,y)}} \implies minus(\omega(x), \omega(y)) \implieas minus(k+1, \omega(y)) \implies k + 1 - \omega(y) \\
& \quad \omega'(b) = I(\delta, \omega, y) \implies \omega(y) \\
= & I(\delta, \omega, \underline{\texttt{plus(a, 1)}}) \\
& \quad \omega''(a) = I(\delta, \omega', \underline{\texttt{div(a,b)}}) \implies div(\omega'(a), \omega'(b)) \implies \omega'(a) \div \omega'(b) \implies (k+1 - \omega(y)) \div \omega'(y)) \\
= & \omega''(a) + 1 \\
= & ((k+1) \div \omega(y)) + 1 \\
= & \frac{k}{\omega(y)} + \frac{1}{\omega(y)} - \frac{\omega(y)}{\omega(y)} + 1 \\
= & \frac{k}{\omega(y)} + \frac{1}{\omega(y)} - 1 + 1 \\
= & \frac{k+1}{\omega(y)} \\
= & \frac{\omega(x)}{\omega(y)} = \omega(x) \div \omega(y)
\end{align*}
\end{enumerate}

\task
\begin{center}
\textbf{1. (sum(as) != 0) == !allZero(as)}
\end{center}

Lemma 1: \texttt{(x == 0 \&\& y == 0) == (x + y == 0)}    \\
Lemma 2: \texttt{(x != 0 || y != 0) == (x + y != 0)} \\
Base: \texttt{as == Nil}
\begin{verbatim}
sum(Nil) != 0
sum(Nil) != sum(Nil)        def. sum. r.l.
false                       logic/arithm
!true                       logic/arithm
!allZero(Nil)               def allZero. r.l.
\end{verbatim}

\textbf{Induction step:} Show the assumption holds for \texttt{a::as}

\begin{verbatim}
sum(a::as) != 0 
a + sum(as) != 0            def sum. l.r.
a != 0 || sum(as) != 0      Lemma 2 r.l.
a != 0 || !allZero(as)      Hypoth. l.r.
!(a == 0 || allZero(as))    De Morgan's law
!allZero(a::as)             def. allZero. r.l
\end{verbatim}

\newpage

\begin{center}
\textbf{2. bigOr(push(as, bs)) == bigOr(as) $||$ bigOr(bs)}
\end{center}

Base: \texttt{as == Nil}
\begin{verbatim}
bigOr(push(Nil, bs))
bigOr(bs)                   def. push. l.r
false   || bogOr(bs)        logic/arithm
bogOr(Nil) || bigOr(bs)     def bigOr r.l
\end{verbatim}

\textbf{Induction step:} Show the assumption holds for \texttt{a::as}

\begin{verbatim}
bigOr(push(a::as, bs))
bigOr(a::push(as, bs))          def. push. l.r.
a || bigOr(push(as, bs))        def. bigOr. l.r.
a || bigOr(as) || bigOr(bs)     Hypoth. l.r.
bigOr(a::as) || bigOr(bs)       def. bigOr. r.l.
\end{verbatim}

\vspace{2cm}

\begin{center}
\textbf{3. bigOr(as) == bigOr(evens(as)) $||$ bigOr(odds(bs))}
\end{center}

Base: \texttt{as == Nil}
\begin{verbatim}
bigOr(Nil)
false                                       def. bigOr. l.r
false   || false                            logic/arithm
bogOr(Nil) || false                         def bigOr r.l
bogOr(Nil) || bigOr(Nil)                    def. bigOr. r. l.
bogOr(evens(Nil)) || bigOr(Nil)             def. evens. r. l.
bogOr(evens(Nil)) || bigOr(odds(Nil))       def. odds. r. l.
\end{verbatim}

\textbf{Induction step:} Show the assumption holds for \texttt{a::y::as}

\begin{verbatim}
bigOr(a::y::as)
a || bigOr(y::as)                                   def bigOr. l.r.
a || y || bigOr(as)                                 def bigOr. l.r.
a || bigOr(evens(as)) || y || bigOr(odds(as))       Hypoth. l.r.
bigOr(a::evens(as)) || y || bigOr(odds(as))         def. evens. r.l.
bigOr(a::evens(as)) || bigOr(y::odds(as))           def. odds. r.l.
bigOr(evens(a::y::as)) || bigOr(y::odds(as))        def. evens r.l.
bigOr(evens(a::y::as)) || bigOr(odds(a::y::as))     def. odds r.l.
\end{verbatim}


\end{document}
